\begin{document}

%------------------------------------------------------------------ 页眉
\resumeheader{宋睿}%
  {13298121974~|~songrui0006@163.com~|~中共党员~|~女~|~2000.06}%
  {photo.jpg}%
  {工艺技术主管}

%------------------------------------------------------------------ 教育背景
\ressection{教育背景}

\entry{中国海洋大学}{2024.09 至 2027.06}{化学工程与技术（硕士）~|~GPA 3.58}
\begin{bullets}
  \item \textbf{主修课程}：分子模拟计算（90）、高等仪器分析（93）、数值分析（96）
\end{bullets}

\entry{河南工业大学}{2018.09 至 2022.07}{化学工程与工艺（本科）~|~GPA 3.76（专业前 10\%）}
\begin{bullets}
  \item \textbf{主修课程}：物理化学（95）、化工原理（91）、高等数学（96）、化工制图（90）、实验设计与数据处理（89）
\end{bullets}

%------------------------------------------------------------------ 科研经历
\ressection{科研经历}

\entry{电催化硝酸根还原合成氨（eNO$_3$RR）}{2025.12 至今}{中国科学院青岛生物能源与过程研究所}
\begin{bullets}
  \item 设计铜基催化剂，通过缺陷与界面工程调控电子结构；系统筛选不同形貌 Cu 的性能并确定适宜反应条件，引入 Co(OH)$_2$ 改性后氨产率提升 7 倍、法拉第效率接近 100\%
  \item 建立质谱原位监测方法追踪反应路径，结合 DFT 计算定位活性位点，形成「合成—测试—机理」的数据闭环
\end{bullets}

\entry{电催化氮气氧化合成硝酸根（eNOR）}{2024.08 至 2025.11}{中国海洋大学}
\begin{bullets}
  \item 制备 $\delta$-MnO$_2$/MoS$_2$ 异质结催化剂，通过组分比例调控筛选最优配比；引入过一硫酸氢钾（PMS）高级氧化工艺提升反应性能
  \item 产率达 198.5 $\mu$g·h$^{-1}$·mg$^{-1}$，法拉第效率 19.31\%；结合原位拉曼与 DFT 计算揭示反应机理
\end{bullets}

%------------------------------------------------------------------ 项目经历
\ressection{项目经历}

\entry{年产 10 万吨异丙醇工艺设计}{2021.02 至 2021.07}{第十五届全国大学生化工设计大赛~|~小组成员}
\begin{bullets}
  \item 经文献调研选定丙酮催化加氢为核心工艺，综合市场、政策与供应链条件确定经济生产规模
  \item 独立完成可行性研究报告、经济评价、「三废一噪」处理方式分析及 HSE 分析全套技术文档，形成完整工业化设计方案；获华中赛区三等奖
\end{bullets}

\entry{海藻酸钠水凝胶的制备及吸附性能研究}{2019.12 至 2022.07}{项目负责人~|~河南工业大学 党丹老师课题组}
\begin{bullets}
  \item 以离子交联法制备水凝胶，系统考察交联条件与吸附环境对性能的影响；优化条件下对刚果红去除率达 95\%，理论最大吸附量 180 mg/g
  \item 结合吸附等温与动力学模型拟合及 ESEM、BET、FTIR 表征，阐明吸附机理
  \item 统筹 10 人以上团队并培养 8 名低年级组成员；发表论文并申请专利
\end{bullets}

%------------------------------------------------------------------ 学生工作
\ressection{学生工作}

\skillrow{研究生兼职辅导员（2025.03 至 2025.08）}%
  {中国海洋大学化学化工学院。参与党员发展全流程，承担会务协助、信息整理、文档校对、考核评奖与档案管理}
\skillrow{党支部组织委员兼纪检委员（2024.08 至今）}%
  {2024 级硕士研究生第三支部。组织主题党日活动，负责党员信息库动态更新与党费收缴公示}

%------------------------------------------------------------------ 专业技能
\ressection{专业技能}

\skillrow{工艺与材料合成}{异质结催化剂合成、缺陷与界面工程调控、反应条件筛选与参数优化、离子交联、自由基聚合}
\skillrow{表征与测试}{循环伏安等电化学测试、原位拉曼、原位质谱、ESEM、BET、FTIR}
\skillrow{计算与数据分析}{DFT 计算、分子模拟、吸附等温与动力学模型拟合、实验设计与数据处理}
\skillrow{工程文档}{独立撰写可行性研究报告、经济评价、HSE 分析等工程技术文档}

%------------------------------------------------------------------ 荣誉与成果
\ressection{荣誉与成果}

\skillrow{学术成果}{期刊论文 1 篇（Adv.\ Mater.\ Sci.\ Eng., 2023）、专利 1 项；硕士期间论文在投、在写各 1 篇}
\skillrow{竞赛与荣誉}{「互联网+」大赛河南赛区一等奖、全国化工设计大赛华中赛区三等奖；优秀研究生、优秀兼职辅导员；校优秀学生奖学金 2 次、优秀共青团员 2 次}

\end{document}
